%% kapa: exam=1 round=32 year=2021 period=2 form=A subject=accounting subject_ko=회계학 seq=20 number=60
%% source: Q-Net 감정평가사 기출문제 (cst00309) artlSeq=5208397 file=extracted/1008/5208397/A형/2021년도 제32회 감정평가사 1차 2교시 A형.pdf
%% conversion: kordoc 4.24.51 -> tools/gen1.py (auto); review=proofread-sonnet; flags=none
\begin{question}{60}
(주)감평은 20×1년 중 연구개발비를 다음과 같이 지출하였다.

{\small\centering\begin{tabular}{|c|c|r|c|}\hline
지출시기 & 구분 & \multicolumn{1}{c|}{금액} & 비고 \\ \hline
1월 초 $\sim$ 6월 말 & 연구단계 & ￦50,000 & \\ \hline
7월 초 $\sim$ 9월 말 & 개발단계 & 100,000 & 자산인식 요건 미충족함 \\ \hline
10월 초 $\sim$ 12월 말 & 개발단계 & 50,000 & 자산인식 요건 충족함 \\ \hline
\end{tabular}\par}
\vspace{4pt}
(주)감평은 20×2년 말까지 ￦100,000을 추가 지출하고 개발을 완료하였다. 무형자산으로 인식한 개발비(내용연수 10년, 잔존가치 ￦0, 정액법 상각)는 20×3년 1월 1일부터 사용이 가능하며, 원가모형을 적용한다. 20×3년 말 현재 개발비가 손상징후를 보였으며 회수가능액은 ￦80,000이다. 20×3년 인식할 개발비 손상차손은?
\begin{choices}
  \item ￦50,000
  \item ￦50,500
  \item ￦53,750
  \item ￦55,000
  \item ￦70,000
\end{choices}
\end{question}
\answer{4}
